\chapter{The Virtual Brain Computing Framework}
\label{ch:virtual_brain}

\papernote{Source paper: \emph{Virtual Brain Computing Framework} --- \texttt{mekaneck/docs/virtual-brain/}}

\begin{chapterabstract}
The Virtual Brain Computing Framework synthesises all preceding results into a complete computational neuroscience platform in which experiment, observation, and computation are mathematically identical. Consciousness is the intersection of multiple neural propagation modality surfaces. Mental states are trajectory-terminus-memory triples. The \Hplus flux at $4.06 \times 10^{13}$~Hz provides the emotional substrate; memory is accumulated emotional change enabling temporal differentiation. The framework includes the Neural Partition Language (NPL) for specifying neural computations, the S-distance neural metric, and multimodal propagation surface intersection algorithms. The mind-body problem is resolved as a corollary of the charge-naming circuit isomorphism.
\end{chapterabstract}

\input{../mekaneck/docs/virtual-brain/sections/consciousness-intersection.tex}
\input{../mekaneck/docs/virtual-brain/sections/trajectory-mental-states.tex}
\input{../mekaneck/docs/virtual-brain/sections/sensory-sufficiency.tex}
\input{../mekaneck/docs/virtual-brain/sections/dreams-and-memory.tex}
\input{../mekaneck/docs/virtual-brain/sections/neural-partition-language-specification.tex}
